\documentclass[11pt,a4paper,english,hyperref]{article}

\newcommand\blattNummer{8}
\newcommand\abgabeDatum{8.3.2019}
\newcommand\ersteAufgabeNummer{1}

%%%%%%%%%%%%%%%%%%%%%%%%%%%%%%%%%%%%%%%%%%%%%%%%%%%%%%%%%%%%%%%%%%%%%%%%%% 
\usepackage{babel}
\usepackage{exscale}
\usepackage{amsmath,amsfonts,amssymb}
\usepackage[thref, amsmath, amsthm, thmmarks]{ntheorem}
\usepackage{a4wide}
\usepackage{epic,eepic}
\usepackage{graphicx} % einbinden von Grafiken mit \includegraphics
\usepackage{url}
\usepackage{mathtools}
\usepackage{etoolbox}
\usepackage[colorlinks]{hyperref}
\usepackage{enumitem}
\usepackage[capitalize]{cleveref}

\pagestyle{empty}
\parindent0em
\parskip1ex
\addtolength{\topmargin}{0cm}
\addtolength{\oddsidemargin}{0cm}
\addtolength{\textheight}{2cm}
\addtolength{\textwidth}{0cm}

% neuer Zähler für die Aufgaben
\newcounter{aufgabeNummer}
\setcounter{aufgabeNummer}{\ersteAufgabeNummer}
\addtocounter{aufgabeNummer}{-1}

% in enumerate-Umgebungen erst Kleinbuchstaben, dann kleine römische Zahlen
\renewcommand{\labelenumi}{(\alph{enumi})\,}
\renewcommand{\labelenumii}{(\roman{enumii})\,}
\crefname{enumerate}{}{}

% Stil für Definitionen, Sätze usw.
\theoremstyle{break}   % Titelzeile abgesetzt
\newtheorem{Satz}{Proposition}
\newtheorem{Lemma}{Lemma}
\newtheorem{Definition}{Definition}
\theorembodyfont{\upshape}  % normaler Font f"ur Beispiel, Aufgabe
\newtheorem{Beisiep}{Example}
\newtheorem{Aufgabe}[aufgabeNummer]{Exercise}

% Mathematische Kürzel
\newcommand\R{\mathbb{R}}
\newcommand\C{\mathbb{C}}
\newcommand\Q{\mathbb{Q}}
\newcommand\N{\mathbb{N}}
\newcommand\Z{\mathbb{Z}}
\renewcommand\P{\mathbb{P}}
\newcommand\T{{\mathcal T} }

\newcommand\eps{{\varepsilon}}                                  % machine precision?
\newcommand\im{{\sl i\;}}                                        % imaginary unit?
\newcommand\e{{\rm e}}                                          % Euler number
\newcommand\D{{\rm D\;}}
\DeclarePairedDelimiter\abs{\lvert}{\rvert}
\DeclarePairedDelimiter\norm{\lVert}{\rVert}
\DeclarePairedDelimiter\dotProductD{\langle}{\rangle}           % Delimiters for dot product

\DeclareMathOperator{\spann}{span}
\DeclareMathOperator{\cond}{cond}
\DeclareMathOperator{\round}{round}
\DeclareMathOperator{\supp}{supp}


\usepackage{scrlayer-scrpage}
\pagestyle{scrheadings}
\clearpairofpagestyles
\ohead{Name: ................................................}


\begin{document}
\parbox{0ex}{ \scalebox{.27}{  \includegraphics[width=.5\textwidth,angle=-90]{cm-logo2-eps-converted-to.pdf} }   }   \\
\parbox{25ex}{
  Prof.~Dr.~S.~Sauter\\
  Institut für Mathematik\\
  Universität Zürich
}
% 
\rule[0cm]{0.cm}{.01cm}                  
\hfill  \parbox{0.6\textwidth}{
  {\sf\LARGE\bfseries Numerik I}
  {\sf\Large\bfseries \;\;--\;\; Midterm}\\[1.5ex]
}
\vspace{5ex}


% \scalebox{.27}{ \put(1250,0){ \includegraphics[width=.5\textwidth,angle=-90]{cm-logo2.eps} }  }


%%%%%%%%%%%%%%%%%%%%%%%%%%%%%%%%%%%%%%%%%%%%%%%%%%%%%%%%%%%%%%%%%%%%%%%%%%% 
% 
% Aufgaben 
% 
%%%%%%%%%%%%%%%%%%%%%%%%%%%%%%%%%%%%%%%%%%%%%%%%%%%%%%%%%%%%%%%%%%%%%%%%%%% 
Exercises marked with $^*$ give extra points.

\begin{Aufgabe}[Computational task]
Let $a$ be a real positive parameter and $f_a:\left[ -1, 1 \right] \to \R$ the given function $f_a(x):= \arctan (a\sin(x)).$ For $n\in \N$ let $p_{a,n}\in \P_n$ be the interpolation polynomial with the following properties
$$
p_{a,n}(x_i)=f_a(x_i)\qquad \text{for nodes }x_i=-1+\frac{2i}{n}\qquad 0\leq i \leq n.
$$
You find the function computing the divided differences and the corresponding interpolation polynomial in your folder (\textbf{Newton\_Interpol.m} or \textbf{Newton\_Interpol.py})
\begin{enumerate}
\item (3 points) Plot the functions $f_a$ and $p_{a,n}$ for the parameter $a=1$ and polynomial degrees $n=4,6,8,10,12$. Store the figure for \(n=8\) as \textbf{fig1a.eps}). Compute the error
$$
E_{a,n}:=\max_{-500\leq j \leq 500} \left| (f_a-p_{a,n})\left(\frac{j}{500}\right) \right|
$$
with respect to \(n\) into a table. Is the decay of the error exponential or polynomial with respect to $n$?
\item (1 point) Estimate, through computer experiments, the approximated value of $a_{max}$, such that $\forall a\leq a_{max}$ the interpolation polynomial $p_{a,n}$ is converging on $\left[ -1,1 \right]$. 
\item(2 points) Fix $a=a_{max}/2$ and find an upper bound for the error  with respect to $n$, estimating $\beta$ and $C$ in the following inequality:
$$
E_{a,n}\leq C e^{-\beta n}=:B(C,\beta,n).
$$
Save a figure with the error and the bound $B(C,\beta,\cdot)$, with respect to $n$, as \textbf{fig1c.eps}.
\end{enumerate}
\end{Aufgabe}
\newpage
\begin{Aufgabe}[Computational task]
Let $$
A_n=(a_{ij})_{i,j=1}^n,\qquad a_{ij}:=\begin{cases} 2 & \text{if }i=j,\\ -1 &\text{if }|i-j|=1, \\ 0 & \text{else}
\end{cases} \qquad \text{ and }b=(b_i)_{i=1}^n, \qquad b_{i}:=\begin{cases} 1 & \text{if }i \text{ even},\\ 0 & \text{if }i\text{ odd.}
\end{cases}
$$
\begin{enumerate}
\item (2 points) Test the running time for the given QR decomposition (\textbf{solveLinSysQR.m} or \textbf{solveLinSysQR.py}) to solve the system \(A_n x = b\) for different values of \(n\) (take \(n \geq 300\)) and verify the asymptotic behaviour of the computational cost with respect to the dimension \(n\) of the system.
\item[(b$^*$)] (6 points) Fix \(n=100\). We consider a perturbed vector on the right hand side
\begin{align*}
b^{\delta}=(b_i)_{i=1}^n, \qquad b_{i}:=\begin{cases} 1 & \text{if }i \text{ even},\\ \delta & \text{if }i\text{ odd}
\end{cases}
\end{align*}
for \(0<\delta< 0.1\). Compute the error \(\|x-x^{\delta}\|\) for different values of \(\delta\), where \(x\) is the solution of \(A_n x= b\), and \(x^{\delta}\) is the solution of \(A_n x^{\delta} = b^{\delta}\). Plot the error with respect to \(\delta\). What can you say about the asymptotic behaviour of the error for \(\delta \rightarrow 0\)? Repeat your experiment with the matrix
$$
\tilde{A}_n=(a_{ij})_{i,j=1}^n,\qquad a_{ij}:=\begin{cases} 4 & \text{if }i=j\\ -1 &\text{if }|i-j|=1 \\ 0 & \text{else}
\end{cases}\
$$
and the same right hand side \(b\) and \(b^{\delta}\). Compare the error \(\|x-x^{\delta}\|\)  and \(\|\tilde{x}-\tilde{x}^{\delta}\|\) in a table. How do you explain the difference of the magnitudes?
\end{enumerate}
\end{Aufgabe}
\newpage 
\begin{Aufgabe}[Mixed task]
Let $I:=[0,1]$. We want to approximate the integral of a given function $f:I\to \mathbb{R}$ 
$$\mathcal{I}(f):=\int_0 ^ 1 \omega(x) f(x)dx
$$
using a quadrature of the form
$$\mathcal{I}(f)\approx Q(f)=\alpha f\left(x_0\right)
$$
\begin{enumerate}
\item (2 points) Assume that we know the values \(\int_0^1 \omega(x)dx\) and \(\int_0^1 x \omega(x)dx\). Find a formula for \(x_0\) and \(\alpha\) such that the quadrature has maximal degree of exactness. Compute \(x_0\) and \(\alpha\) explicitly for the choice \(\omega(x) = 1+x\).
\end{enumerate}
\noindent From now on, we consider \(\omega(x) = 1\).
\begin{enumerate}
\setcounter{enumi}{1}
\item (2 points) Write a program that computes the composite quadrature \(Q^n(f)\) (based on \(Q(f)\)) on \(n\) equidistant subintervals. 
%INPUT=\{$n$ the number of subintervals, the anonymous function (MATLAB) or the lambda function (Pyhton)\footnote{ex. MATLAB: \texttt{f=@(x) x.\^{}2 }\\ ex. Python : \texttt{f=lambda x: x**2} } $f : I\to \mathbb{R}$  \}; OUTPUT=\{$Q_M^n(f)$\}.
\item (2 points) Run the program for function $f(x):=\dfrac{1}{1 + x^2}$. The exact value of \(\mathcal{I}\) is given by $\mathcal{I}(f)=\pi/4$. Compute error $e^n(f):=|Q^n(f)-\mathcal{I}(f)|$ for different values of $n$. What can you say about the convergence of the error with respect to $h$, where $h$ is the length of each subinterval?
\item (2 points) (Theoretical task) Prove that for $f\in \mathcal{C}^\infty (I)$ $$e^n(f)\leq Ch^2.$$ %
\item (2 points) Use the quadrature for the function $f(x):=x^{2/3}$. Compute analytically the exact value of $I(f)$. Estimate numerically the exponent $r$ in the following inequality 
$$e^n(f) \leq Ch^r.$$
Explain why $r<2$.
\end{enumerate}
\end{Aufgabe}


\end{document}